\documentclass[11pt]{article}
\usepackage[margin=1in]{geometry}
\usepackage{amsmath,amssymb,amsthm,mathtools}
\usepackage{bm}
\usepackage{enumitem}
\usepackage[colorlinks=true,citecolor=blue,urlcolor=blue,linkcolor=blue]{hyperref}

%%% ---- macros ----
\newcommand{\R}{\mathbb{R}}
\newcommand{\E}{\mathbb{E}}
\newcommand{\Var}{\operatorname{Var}}
\newcommand{\Cov}{\operatorname{Cov}}
\newcommand{\tr}{\operatorname{tr}}
\newcommand{\diag}{\operatorname{diag}}
\newcommand{\dTV}{d_{\mathrm{TV}}}
\newcommand{\op}{\mathrm{op}}
\newcommand{\X}{\bm{X}}
\newcommand{\Z}{\bm{Z}}
\newcommand{\y}{\bm{y}}
\newcommand{\bSig}{\bm{\Sigma}}
\newcommand{\hSig}{\widehat{\bm{\Sigma}}}
\newcommand{\bal}{\bm{\alpha}}
\newcommand{\beps}{\bm{\varepsilon}}
\newcommand{\bone}{\bm{1}}
\newcommand{\bx}{\bm{x}}
\newcommand{\bz}{\bm{z}}
\newcommand{\be}{\bm{e}}
\newcommand{\bv}{\bm{v}}
\newcommand{\br}{\bm{r}}
\newcommand{\bc}{\bm{c}}
\newcommand{\bxi}{\bm{\xi}}
\newcommand{\normsq}[1]{\|#1\|^{2}}

\newtheorem{theorem}{Theorem}
\newtheorem{lemma}[theorem]{Lemma}
\newtheorem{proposition}[theorem]{Proposition}
\newtheorem{corollary}[theorem]{Corollary}
\newtheorem{remark}[theorem]{Remark}
\newtheorem{assumption}{Assumption}[section]

\title{\textbf{Asymptotic Normality of the Centered LDSC Estimator}\\[6pt]
\large Under Fixed Genetic Effects and Block-Diagonal Linkage Disequilibrium}
\author{}
\date{}

\begin{document}
\maketitle

\begin{abstract}
We establish the asymptotic normality of the centered linkage disequilibrium
score regression (LDSC) estimator of heritability under a high-dimensional
linear model with \emph{fixed} genetic effects and a block-diagonal population
covariance structure.  The proof proceeds in three stages: (i)~computation of
the first two unconditional moments of the squared marginal association
statistics using Wishart moment formulas (Letac \& Massam, 2004;
Dicker, 2014); (ii)~a martingale central limit theorem that exploits the block
independence of the genotype matrix; (iii)~a quantitative normal
approximation via Chatterjee's (2009) second-order Poincar\'e inequality.
Both asymptotic normality arguments are presented in full detail.
\end{abstract}

\tableofcontents

%=============================================================================
\section{Model and notation}\label{sec:setup}
%=============================================================================

\subsection{Data generating process}\label{ssec:dgp}

Let $p\ge1$ denote the number of genetic variants (single nucleotide
polymorphisms) and let $n\ge1$ and $n_r\ge1$ denote, respectively, the sizes of
the genome-wide association study (GWAS) cohort and an independent reference
panel.  We observe the following data.
\begin{enumerate}[label=(\roman*)]
\item \emph{GWAS genotype matrix.}\; $\X\in\R^{n\times p}$ with independent
  rows $\bx_1,\dots,\bx_n\overset{\mathrm{iid}}{\sim}N(\bm{0},\bSig)$, where
  $\bSig\in\R^{p\times p}$ is a positive-definite covariance matrix.

\item \emph{Phenotype vector.}\; $\y=\X\bal+\beps\in\R^{n}$, where
  $\bal=(\alpha_1,\dots,\alpha_p)^{T}\in\R^p$ is a \emph{fixed}
  (non-random) vector of genetic effects and
  $\beps=(\varepsilon_1,\dots,\varepsilon_n)^{T}\sim N(\bm{0},\sigma_\varepsilon^2 I_n)$
  is independent of $\X$.

\item \emph{Reference panel.}\; $\Z\in\R^{n_r\times p}$ with independent rows
  $\bz_1,\dots,\bz_{n_r}\overset{\mathrm{iid}}{\sim}N(\bm{0},\bSig)$,
  independent of~$(\X,\beps)$.
\end{enumerate}

\subsection{Block-diagonal covariance structure}\label{ssec:block}

Assume that the $p$ variants can be partitioned into $B$ non-overlapping blocks
$S_1,\dots,S_B$ with $|S_b|=p_b$ and $\sum_{b=1}^{B}p_b=p$, and that
\begin{equation}\label{eq:blockSig}
  \bSig = \diag(\bSig_1,\dots,\bSig_B),
\end{equation}
where $\bSig_b\in\R^{p_b\times p_b}$ is the within-block covariance matrix.
The block-diagonal structure implies that
$\X=(\X^{(1)},\dots,\X^{(B)})$ with $\X^{(b)}\in\R^{n\times p_b}$
mutually independent, and similarly for~$\Z$.  We write
$\bal_{S_b}\in\R^{p_b}$ for the sub-vector of effects in block~$b$.

\subsection{Observed statistics}\label{ssec:stats}

Define the reference-panel sample covariance matrix
\begin{equation}\label{eq:hSig}
  \hSig = \frac{1}{n_r}\,\Z^{T}\Z \;\in\R^{p\times p},
\end{equation}
and for each variant $j=1,\dots,p$:
\begin{align}
  \hat{a}_j &= \frac{1}{n}\,\X_j^{T}\y
    &&\text{(marginal association statistic)}, \label{eq:aj}\\[4pt]
  \hat\ell_j &= \sum_{i=1}^{p}\hSig_{ji}^{2}
    &&\text{(estimated LD score)}, \label{eq:lj}
\end{align}
where $\X_j$ denotes the $j$-th column of~$\X$.

Collect these into vectors
$\hat{\bm{w}}=(\hat{a}_1^{2},\dots,\hat{a}_p^{2})^{T}$,
$\hat{\bm\ell}=(\hat\ell_1,\dots,\hat\ell_p)^{T}$,
and define the cross-sectional means
\[
  \mu_{\hat\ell}=\frac{1}{p}\,\bone_p^{T}\hat{\bm\ell},
  \qquad
  \hat\mu_w=\frac{1}{p}\,\bone_p^{T}\hat{\bm{w}}.
\]

\begin{equation}\label{eq:LDSC}
\boxed{\;
  \hat\sigma_\alpha^2
  \;=\;\frac{(\hat{\bm\ell}-\mu_{\hat\ell}\bone)^{T}
             (\hat{\bm{w}}-\hat\mu_w\bone)}
            {(\hat{\bm\ell}-\mu_{\hat\ell}\bone)^{T}
             (\hat{\bm\ell}-\mu_{\hat\ell}\bone)}
\;}
\end{equation}
is the \emph{centered LDSC estimator}.

\subsection{Additional notation}\label{ssec:notation}

We employ the following shorthand throughout the paper.

\medskip
\noindent\textbf{Population quantities.}\;
\begin{alignat}{3}
  a_j^{*} &= (\bSig\bal)_j
    = \textstyle\sum_{k=1}^{p}\Sigma_{jk}\alpha_k,
    &\qquad& \ell_j &&= \textstyle\sum_{i=1}^{p}\Sigma_{ji}^{2}
      \;\;\text{(population LD score)}, \label{eq:pop_LD}\\[2pt]
  \tau^{2} &= \bal^{T}\bSig\bal,
    && \sigma_y^{2} &&= \tau^{2}+\sigma_\varepsilon^{2},\label{eq:tau}\\[2pt]
  \bar\ell &= \textstyle\frac{1}{p}\sum_{j=1}^{p}\ell_j,
    && V_\ell &&= \textstyle\frac{1}{p}\sum_{j=1}^{p}(\ell_j-\bar\ell)^{2}.
    \label{eq:Vell}
\end{alignat}

\medskip
\noindent\textbf{Centered LD scores and quadratic form.}\;
Set $c_j=\hat\ell_j-\mu_{\hat\ell}$ and $S_c=\sum_{j=1}^{p}c_j^{2}=\|\bc\|^{2}$.
Since $\sum_{j=1}^{p}c_j=0$, the estimator~\eqref{eq:LDSC} can be written as
\begin{equation}\label{eq:LDSC2}
  \hat\sigma_\alpha^2
  = \frac{1}{S_c}\sum_{j=1}^{p}c_j\,\hat{a}_j^{2}.
\end{equation}
Introducing $D=\diag(c_1,\dots,c_p)$, we have
\begin{equation}\label{eq:quadform}
  \sum_{j=1}^{p}c_j\,\hat{a}_j^{2}
  = \frac{1}{n^{2}}\,\y^{T}\X D\X^{T}\y.
\end{equation}

\medskip
\noindent\textbf{Wishart matrices.}\;
$W=\X^{T}\X\sim\text{Wishart}_p(n,\bSig)$ and $W_Z=\Z^{T}\Z\sim\text{Wishart}_p(n_r,\bSig)$, independent.

%=============================================================================
\section{Regularity assumptions}\label{sec:assumptions}
%=============================================================================

\begin{assumption}[Bounded block sizes]\label{ass:block}
There exists a constant $p_{\max}<\infty$ such that $\sup_{b}\,p_b\le p_{\max}$
for all~$p$.  In particular, $B=\Theta(p)$ as $p\to\infty$.
\end{assumption}

\begin{assumption}[Dimensional regime]\label{ass:regime}
As $n,p\to\infty$, we have
$\displaystyle\frac{p}{n}\to\gamma\in(0,\infty)$
and
$\displaystyle\frac{p}{n_r}\to\gamma_r\in(0,\infty)$.
\end{assumption}

\begin{assumption}[Spectral regularity]\label{ass:spectral}
There exist constants $0<\underline\lambda\le\bar\lambda<\infty$ such that for
every block~$b$ and all~$p$,
\[
  \underline\lambda\le\lambda_{\min}(\bSig_b)
  \le\lambda_{\max}(\bSig_b)\le\bar\lambda.
\]
\end{assumption}

\begin{assumption}[Effect-size regularity]\label{ass:effects}
The empirical second moment
$\sigma_\alpha^{2,\circ}=p^{-1}\sum_{j=1}^{p}\alpha_j^{2}$ converges to a
finite limit as $p\to\infty$.  More generally, for every block-level function
$f:\R^{p_{\max}}\times\R^{p_{\max}\times p_{\max}}\to\R$ that is bounded
and continuous, the limit
\[
  \frac{1}{B}\sum_{b=1}^{B}f(\bal_{S_b},\bSig_b)
  \;\longrightarrow\; \bar{f}
\]
exists as $p\to\infty$.
\end{assumption}

\begin{assumption}[Non-degeneracy of LD-score variation]\label{ass:ldvar}
There exists a constant $c_\ell>0$ such that $V_\ell\ge c_\ell$ for all
large~$p$.
\end{assumption}

%=============================================================================
\section{Expected value of the LDSC estimator}\label{sec:mean}
%=============================================================================

\begin{lemma}[Mean of $\hat{a}_j^{2}$]\label{lem:Eaj2}
Under the model of Section~\ref{sec:setup},
\begin{equation}\label{eq:Eaj2}
  \E[\hat{a}_j^{2}]
  = \Bigl(1+\frac{1}{n}\Bigr)(a_j^{*})^{2}
    + \frac{\Sigma_{jj}}{n}\,\sigma_y^{2}.
\end{equation}
\end{lemma}

\begin{proof}
Since $\hat{a}_j = n^{-1}\X_j^{T}(\X\bal+\beps)$, the conditional moments
given~$\X$ are
\begin{equation}\label{eq:cond_moments_aj}
  \E[\hat{a}_j\mid\X]
  = \frac{1}{n}\X_j^{T}\X\bal
  = \sum_{k=1}^{p}\alpha_k\,\hat\Sigma_{jk}^{(X)},
  \qquad
  \Var(\hat{a}_j\mid\X)
  = \frac{\sigma_\varepsilon^{2}\|\X_j\|^{2}}{n^{2}},
\end{equation}
where $\hat\Sigma_{jk}^{(X)}=n^{-1}\X_j^{T}\X_k$.  By the law of total
expectation,
\begin{equation}\label{eq:total_var}
  \E[\hat{a}_j^{2}]
  = \underbrace{\E\bigl[\Var(\hat{a}_j\mid\X)\bigr]}_{(I)}
    + \underbrace{\E\bigl[\bigl(\E[\hat{a}_j\mid\X]\bigr)^{2}\bigr]}_{(II)}.
\end{equation}

\medskip\noindent\emph{Term~(I).}\;
Since $\E[\|\X_j\|^{2}]=n\Sigma_{jj}$, we have
\begin{equation}\label{eq:term1}
  (I) = \frac{\sigma_\varepsilon^{2}}{n^{2}}\,\E[\|\X_j\|^{2}]
       = \frac{\sigma_\varepsilon^{2}\Sigma_{jj}}{n}.
\end{equation}

\medskip\noindent\emph{Term~(II).}\;
We need $\E[W_{jk}W_{jl}]$ for $W=\X^{T}\X\sim\text{Wishart}_p(n,\bSig)$.
Using the fundamental Wishart moment formula~(S16) of Dicker~(2014), which
states that for symmetric matrices $H_1,H_2$,
\begin{equation}\label{eq:S16_basic}
  \E[\tr(WH_1)\tr(WH_2)]
  = 2n\,\tr(\bSig H_1\bSig H_2)+n^{2}\tr(\bSig H_1)\tr(\bSig H_2),
\end{equation}
we set $H_1=\be_j\be_k^{T}+\be_k\be_j^{T}$ and $H_2=\be_j\be_l^{T}+\be_l\be_j^{T}$
(where $\be_j$ is the $j$-th standard basis vector in~$\R^p$) to derive
\begin{equation}\label{eq:wishart_second}
  \E[W_{jk}W_{jl}]
  = n^{2}\Sigma_{jk}\Sigma_{jl}+n\bigl(\Sigma_{jj}\Sigma_{kl}+\Sigma_{jk}\Sigma_{jl}\bigr).
\end{equation}
Consequently,
\begin{align}
  (II) &= \frac{1}{n^{2}}\sum_{k,l=1}^{p}\alpha_k\alpha_l\,\E[W_{jk}W_{jl}]
  \notag\\
  &= \sum_{k,l}\alpha_k\alpha_l\Sigma_{jk}\Sigma_{jl}
     + \frac{1}{n}\Bigl(\Sigma_{jj}\sum_{k,l}\alpha_k\alpha_l\Sigma_{kl}
     + \sum_{k,l}\alpha_k\alpha_l\Sigma_{jk}\Sigma_{jl}\Bigr)
  \notag\\
  &= (a_j^{*})^{2} + \frac{1}{n}\bigl[\Sigma_{jj}\tau^{2}+(a_j^{*})^{2}\bigr].
  \label{eq:term2}
\end{align}
Adding~\eqref{eq:term1} and~\eqref{eq:term2} and using
$\sigma_y^{2}=\tau^{2}+\sigma_\varepsilon^{2}$ yields~\eqref{eq:Eaj2}.
\end{proof}

\begin{remark}[Interpretation]\label{rem:interpretation}
Equation~\eqref{eq:Eaj2} decomposes as
\[
  \E[\hat{a}_j^{2}]
  = (a_j^{*})^{2}
    + \frac{1}{n}\bigl[(a_j^{*})^{2}+\Sigma_{jj}\sigma_y^{2}\bigr].
\]
The first term is the squared ``true marginal covariance'' between the $j$-th
variant and the phenotype.  The $O(1/n)$ correction is a finite-sample bias
arising from sampling variability in the genotypes and the noise~$\beps$.
\end{remark}

\medskip

Define the \emph{pseudo-true value} of the LDSC estimator by
\begin{equation}\label{eq:sigma0}
  \sigma_{\alpha,0}^{2}
  = \frac{1}{S_c}\sum_{j=1}^{p}c_j\,\E[\hat{a}_j^{2}]
  = \frac{1}{S_c}\Biggl[
      \Bigl(1+\frac{1}{n}\Bigr)\sum_{j=1}^{p}c_j(a_j^{*})^{2}
      + \frac{\sigma_y^{2}}{n}\sum_{j=1}^{p}c_j\Sigma_{jj}
    \Biggr],
\end{equation}
where we have used $\sum_{j}c_j=0$ to eliminate the constant
$\sigma_y^{2}\Sigma_{jj}/n$ from $\E[\hat{a}_j^2]$ that does not interact with
the centering.  Note that $\sigma_{\alpha,0}^{2}$ depends on $\Z$ through~$\bc$.

%=============================================================================
\section{Covariance structure and variance}\label{sec:variance}
%=============================================================================

\subsection{Covariance of squared association statistics}

\begin{lemma}[Covariance of $\hat{a}_j^{2}$ and
$\hat{a}_k^{2}$]\label{lem:cov}
For any $j,k\in\{1,\dots,p\}$,
\begin{align}
  \Cov(\hat{a}_j^{2},\hat{a}_k^{2})
  &= \frac{2}{n}\bigl(\Sigma_{jk}\sigma_y^{2}+n\,a_j^{*}a_k^{*}\bigr)^{2}
     \cdot\frac{1}{n^{2}}
     + O(n^{-3})
  \notag\\[4pt]
  &= \frac{2}{n}\Bigl(\frac{\Sigma_{jk}\sigma_y^{2}}{n}+a_j^{*}a_k^{*}\Bigr)^{\!2}
     + O(n^{-3}).
  \label{eq:cov_leading}
\end{align}
In particular, if $j\in S_b$ and $k\in S_{b'}$ with $b\ne b'$, then
$\Sigma_{jk}=0$ and
\begin{equation}\label{eq:cov_cross}
  \Cov(\hat{a}_j^{2},\hat{a}_k^{2})
  = \frac{2}{n}(a_j^{*})^{2}(a_k^{*})^{2}+O(n^{-3}).
\end{equation}
\end{lemma}

\begin{proof}
We use the law of total covariance:
\begin{equation}\label{eq:total_cov}
  \Cov(\hat{a}_j^{2},\hat{a}_k^{2})
  = \E\bigl[\Cov(\hat{a}_j^{2},\hat{a}_k^{2}\mid\X)\bigr]
    + \Cov\bigl(\E[\hat{a}_j^{2}\mid\X],\,\E[\hat{a}_k^{2}\mid\X]\bigr).
\end{equation}

\medskip
\noindent\textbf{Inner covariance.}\;
Given~$\X$, we have
$\hat{a}_j=n^{-1}\X_j^{T}\y$ with
$\y\mid\X\sim N(\X\bal,\sigma_\varepsilon^{2}I_n)$, so
$(\hat{a}_j,\hat{a}_k)\mid\X$ is jointly normal with
\begin{equation}\label{eq:cond_cov_jk}
  \Cov(\hat{a}_j,\hat{a}_k\mid\X)
  = \frac{\sigma_\varepsilon^{2}}{n^{2}}\,\X_j^{T}\X_k
  = \frac{\sigma_\varepsilon^{2}}{n}\,\hat\Sigma_{jk}^{(X)},
\end{equation}
and conditional means
$\mu_j(\X)=\E[\hat{a}_j\mid\X]=\sum_l\alpha_l\hat\Sigma_{jl}^{(X)}$.

For jointly normal random variables $(U,V)$ with mean $(\mu_U,\mu_V)$ and
covariance $\sigma_{UV}$, the Isserlis--Wick theorem gives
\begin{equation}\label{eq:isserlis}
  \Cov(U^{2},V^{2})
  = 2\sigma_{UV}^{2} + 4\mu_U\mu_V\sigma_{UV}.
\end{equation}
Applying~\eqref{eq:isserlis} conditionally yields
\begin{align}
  \Cov(\hat{a}_j^{2},\hat{a}_k^{2}\mid\X)
  &= 2\Bigl(\frac{\sigma_\varepsilon^{2}}{n}\,\hat\Sigma_{jk}^{(X)}\Bigr)^{\!2}
     + 4\mu_j(\X)\mu_k(\X)\,\frac{\sigma_\varepsilon^{2}}{n}\,\hat\Sigma_{jk}^{(X)}.
  \label{eq:cond_cov_sq}
\end{align}
Taking the expectation over~$\X$ and applying the Wishart moment formula~\eqref{eq:wishart_second} together with the higher-order formulas of
Proposition~S1 of Dicker~(2014) (specifically equations (S17)--(S22), applied
block-by-block using the independence of the blocks), we obtain
\begin{equation}\label{eq:inner_cov}
  \E\bigl[\Cov(\hat{a}_j^{2},\hat{a}_k^{2}\mid\X)\bigr]
  = \frac{2\sigma_\varepsilon^{4}}{n^{2}}
    \Bigl(\Sigma_{jk}^{2}+\frac{\Sigma_{jj}\Sigma_{kk}+\Sigma_{jk}^{2}}{n}\Bigr)
    + \frac{4\sigma_\varepsilon^{2}}{n}\cdot\frac{\Sigma_{jk}}{n}
      \Bigl[a_j^{*}a_k^{*}+\frac{\Sigma_{jk}\tau^{2}+\Sigma_{jj}a_k^{*}\alpha_j
        +\Sigma_{kk}a_j^{*}\alpha_k}{n}\Bigr]
    + O(n^{-3}).
\end{equation}
Here we used the following consequence of~\eqref{eq:wishart_second}:
\[
  \E\bigl[\hat\Sigma_{jk}^{(X)}\hat\Sigma_{jl}^{(X)}\bigr]
  = \Sigma_{jk}\Sigma_{jl}+\frac{\Sigma_{jj}\Sigma_{kl}+\Sigma_{jk}\Sigma_{jl}}{n},
\]
and the third-order Wishart formula (equation~(S22) of Dicker, 2014):
\[
  \E[\tr(W)\,\beta^{T}W\beta]
  = dn^{2}m_1\tau_1^{2}+2n\tau_2^{2},
\]
applied to the relevant sub-blocks, where the notation is as in Proposition~S1
of Dicker~(2014).

\medskip
\noindent\textbf{Outer covariance.}\;
By~\eqref{eq:cond_moments_aj}, the conditional expectation of $\hat{a}_j^{2}$ given $\X$ is
\[
  \E[\hat{a}_j^{2}\mid\X]
  = \Var(\hat{a}_j\mid\X)+\bigl(\E[\hat{a}_j\mid\X]\bigr)^{2}
  = \frac{\sigma_\varepsilon^{2}}{n}\,\hat\Sigma_{jj}^{(X)}
    + \Bigl(\sum_l\alpha_l\hat\Sigma_{jl}^{(X)}\Bigr)^{\!2}.
\]
The outer covariance is therefore
\begin{align}
  &\Cov\bigl(\E[\hat{a}_j^{2}\mid\X],\E[\hat{a}_k^{2}\mid\X]\bigr)
  \notag\\
  &\qquad= \frac{\sigma_\varepsilon^{4}}{n^{2}}\Cov\bigl(\hat\Sigma_{jj}^{(X)},
    \hat\Sigma_{kk}^{(X)}\bigr)
  + \frac{2\sigma_\varepsilon^{2}}{n}
    \Cov\Bigl(\hat\Sigma_{jj}^{(X)},
    \Bigl(\textstyle\sum_l\alpha_l\hat\Sigma_{kl}^{(X)}\Bigr)^{\!2}\Bigr)
  \notag\\[4pt]
  &\qquad\quad
  + \Cov\Bigl(\Bigl(\textstyle\sum_l\alpha_l\hat\Sigma_{jl}^{(X)}\Bigr)^{\!2},
    \Bigl(\textstyle\sum_m\alpha_m\hat\Sigma_{km}^{(X)}\Bigr)^{\!2}\Bigr).
  \label{eq:outer_cov}
\end{align}
Each term in~\eqref{eq:outer_cov} involves the second and fourth moments of
the Wishart matrix~$W$.  Applying Proposition~S1 of Dicker~(2014):

\begin{itemize}
\item For the first term, using~\eqref{eq:wishart_second},
  $\Cov(\hat\Sigma_{jj}^{(X)},\hat\Sigma_{kk}^{(X)})
  = n^{-1}(2\Sigma_{jk}^{2}+O(n^{-1}))$.

\item For the cross-term, using the third-order moment
  formulas~(S22)--(S23) of Dicker~(2014) and the block-diagonal structure,
  one obtains contributions of order~$O(n^{-1})$.

\item For the last term (covariance of two quadratic forms in $W$), the
  Isserlis--Wick theorem for Wishart matrices gives the leading behaviour:
  \begin{equation}\label{eq:outer_leading}
    \Cov\Bigl(\Bigl(\textstyle\sum_l\alpha_l\hat\Sigma_{jl}^{(X)}\Bigr)^{\!2},
    \Bigl(\textstyle\sum_m\alpha_m\hat\Sigma_{km}^{(X)}\Bigr)^{\!2}\Bigr)
    = \frac{2}{n}\bigl[\Sigma_{jk}\tau^{2}+a_j^{*}a_k^{*}\bigr]^{2}
      + O(n^{-2}).
  \end{equation}
  To establish~\eqref{eq:outer_leading}, we use the fact that for
  $Q_1=\bal^{T}\be_jW\be_j^{T}\bSig\bal/n$ and
  $Q_2=\bal^{T}\be_kW\be_k^{T}\bSig\bal/n$
  (shorthand for the relevant quadratic forms),
  the fourth moment formula~(S26) of Dicker~(2014) and the expectation
  formulas~(S20)--(S21) provide the exact expressions needed.
\end{itemize}

\medskip
Combining~\eqref{eq:inner_cov} with~\eqref{eq:outer_leading} and the
remaining terms from~\eqref{eq:outer_cov}, collecting the leading-order
contributions, we arrive at
\begin{equation}\label{eq:cov_full}
  \Cov(\hat{a}_j^{2},\hat{a}_k^{2})
  = \frac{2\sigma_\varepsilon^{4}\Sigma_{jk}^{2}}{n^{2}}
    + \frac{4\sigma_\varepsilon^{2}\Sigma_{jk}\,a_j^{*}a_k^{*}}{n^{2}}
    + \frac{2}{n}\bigl[\Sigma_{jk}\tau^{2}+a_j^{*}a_k^{*}\bigr]^{2}
    + O(n^{-3}).
\end{equation}
Rearranging the first three terms into a single square gives
\begin{align}
  \frac{2\sigma_\varepsilon^{4}\Sigma_{jk}^{2}}{n^{2}}
  + \frac{4\sigma_\varepsilon^{2}\Sigma_{jk}a_j^{*}a_k^{*}}{n^{2}}
  + \frac{2}{n}(a_j^{*}a_k^{*})^{2}
  + \frac{2}{n}\bigl[\Sigma_{jk}^{2}\tau^{4}+2\Sigma_{jk}\tau^{2}a_j^{*}a_k^{*}\bigr]
  \notag\\
  = \frac{2}{n}\Bigl(\frac{\Sigma_{jk}\sigma_\varepsilon^{2}}{n}
    + a_j^{*}a_k^{*}+\Sigma_{jk}\tau^{2}\Bigr)^{2}
  - \frac{2\Sigma_{jk}^{2}\sigma_\varepsilon^{4}}{n^{3}}
  + O(n^{-3}).
  \label{eq:rearrange}
\end{align}
Since $\sigma_y^{2}=\tau^{2}+\sigma_\varepsilon^{2}$, we can write
$\Sigma_{jk}\sigma_\varepsilon^{2}/n+\Sigma_{jk}\tau^{2}
=\Sigma_{jk}(\sigma_y^{2}-\tau^{2})/n+\Sigma_{jk}\tau^{2}
=\Sigma_{jk}\sigma_y^{2}/n+\Sigma_{jk}\tau^{2}(1-1/n)$.  For the leading-order
formula, we absorb the $O(1/n)$ corrections into the remainder, yielding~\eqref{eq:cov_leading}.  When $\Sigma_{jk}=0$ (cross-block), only the
$(a_j^{*}a_k^{*})^{2}$ term survives, giving~\eqref{eq:cov_cross}.
\end{proof}

\subsection{Variance of the LDSC estimator}

By~\eqref{eq:LDSC2} and~\eqref{eq:cov_leading},
\begin{equation}\label{eq:var_hat_sigma}
  \Var(\hat\sigma_\alpha^{2}\mid\Z)
  = \frac{1}{S_c^{2}}\sum_{j,k=1}^{p}c_jc_k\,\Cov(\hat{a}_j^{2},\hat{a}_k^{2}).
\end{equation}

Define the \emph{asymptotic variance parameter}
\begin{equation}\label{eq:Vdef}
  V = \lim_{p\to\infty}\;
  \frac{2p}{V_\ell^{2}}\cdot\frac{1}{p^{2}}
  \sum_{j,k=1}^{p}
    (\ell_j-\bar\ell)(\ell_k-\bar\ell)
    \Bigl(\frac{\Sigma_{jk}\sigma_y^{2}}{n}+a_j^{*}a_k^{*}\Bigr)^{\!2},
\end{equation}
which is well-defined and strictly positive under
Assumptions~\ref{ass:block}--\ref{ass:ldvar}.

\begin{proposition}[Variance convergence]\label{prop:variance}
Under Assumptions~\ref{ass:block}--\ref{ass:ldvar},
\[
  p\,\Var(\hat\sigma_\alpha^{2}\mid\Z)
  \;\xrightarrow{\;\;P\;\;}\; V
  \quad\text{as }n,p\to\infty,
\]
where the convergence in probability is over the randomness of~$\Z$.
\end{proposition}

\begin{proof}
From~\eqref{eq:var_hat_sigma} and Lemma~\ref{lem:cov},
\begin{equation}\label{eq:var_expansion}
  p\,\Var(\hat\sigma_\alpha^{2}\mid\Z)
  = \frac{2p}{n\,S_c^{2}}
    \sum_{j,k=1}^{p}c_jc_k
    \Bigl(\frac{\Sigma_{jk}\sigma_y^{2}}{n}+a_j^{*}a_k^{*}\Bigr)^{\!2}
    + O\Bigl(\frac{p}{n^{2}}\Bigr).
\end{equation}
The remainder $O(p/n^{2})=O(n^{-1})\to0$ by $p/n\to\gamma$.

\medskip
\noindent\textbf{Concentration of $S_c$ and $c_j$.}\;
Since $\hSig$ inherits the block structure from~$\bSig$, the estimated LD score
for variant~$j$ in block~$b$ is
\[
  \hat\ell_j
  = \sum_{i\in S_b}\hSig_{ji}^{2}
  = \ell_j + O_P(n_r^{-1/2}),
\]
where the error rate follows from the Wishart concentration
(Proposition~S1 of Dicker, 2014, applied to the block~$b$ with
$d=p_b\le p_{\max}$, so that all moments of the entries of $\hSig$ are
uniformly controlled).  Crucially, the cross-block entries of~$\hSig$ vanish
identically because $W_Z=\Z^{T}\Z$ inherits the block-diagonal structure.
Hence $c_j=\hat\ell_j-\mu_{\hat\ell}=(\ell_j-\bar\ell)+O_P(n_r^{-1/2})$
uniformly over~$j$, and
\[
  \frac{S_c}{p}
  = \frac{1}{p}\sum_{j=1}^{p}c_j^{2}
  = V_\ell + O_P(n_r^{-1/2})
  \;\xrightarrow{\;P\;}\; V_\ell > c_\ell > 0.
\]
Substituting $c_j\approx\ell_j-\bar\ell$ and $S_c\approx pV_\ell$
into~\eqref{eq:var_expansion}, and using $p/n\to\gamma$, we obtain
convergence to the deterministic limit~$V$ defined in~\eqref{eq:Vdef}.
\end{proof}

%=============================================================================
\section{Asymptotic normality: statement}\label{sec:statement}
%=============================================================================

\begin{theorem}[Asymptotic normality of the LDSC
estimator]\label{thm:main}
Let Assumptions~\ref{ass:block}--\ref{ass:ldvar} hold and let $n,p\to\infty$
with $p/n\to\gamma\in(0,\infty)$ and $p/n_r\to\gamma_r\in(0,\infty)$.  Then
\begin{equation}\label{eq:CLT}
  \sqrt{p}\,\bigl(\hat\sigma_\alpha^{2}-\sigma_{\alpha,0}^{2}\bigr)
  \;\xrightarrow{\;\;d\;\;}\; N(0,\,V),
\end{equation}
where $\sigma_{\alpha,0}^{2}$ is defined in~\eqref{eq:sigma0} and $V$ in~\eqref{eq:Vdef}.
\end{theorem}

We give two proofs: the first uses a martingale central limit theorem and the
second uses Chatterjee's (2009) second-order Poincar\'e inequality.

%=============================================================================
\section{First proof: martingale CLT}\label{sec:mCLT}
%=============================================================================

\subsection*{Step 1: Conditioning on the reference panel}

Since $\Z$ is independent of $(\X,\beps)$, we condition on $\Z$ and treat
$\bc=(c_1,\dots,c_p)^{T}$ as deterministic.  Define the centred and scaled
statistic
\begin{equation}\label{eq:Tp}
  T_p
  = \sqrt{p}\,\bigl(\hat\sigma_\alpha^{2}-\sigma_{\alpha,0}^{2}\bigr)
  = \frac{\sqrt{p}}{S_c}\sum_{j=1}^{p}c_j
    \bigl(\hat{a}_j^{2}-\E[\hat{a}_j^{2}]\bigr).
\end{equation}
It suffices to show $T_p\xrightarrow{d}N(0,V)$ conditionally on~$\Z$,
for almost every realization of~$\Z$.  By Proposition~\ref{prop:variance},
$\Var(T_p\mid\Z)\to V$ in probability.

\subsection*{Step 2: Block decomposition}

For each block $b=1,\dots,B$, define
\begin{equation}\label{eq:Yb}
  Y_b
  = \frac{\sqrt{p}}{S_c}\sum_{j\in S_b}c_j
    \bigl(\hat{a}_j^{2}-\E[\hat{a}_j^{2}]\bigr),
\end{equation}
so that $T_p=\sum_{b=1}^{B}Y_b$ and $\E[Y_b]=0$.

\medskip
\noindent\textbf{Decomposition of $\hat{a}_j$ for $j\in S_b$.}\;
Write $\y=\X^{(b)}\bal_{S_b}+\br^{(b)}$, where
\begin{equation}\label{eq:rb}
  \br^{(b)}
  = \sum_{\substack{b'=1\\b'\ne b}}^{B}\X^{(b')}\bal_{S_{b'}}+\beps
\end{equation}
is independent of $\X^{(b)}$ (since $\X^{(1)},\dots,\X^{(B)},\beps$ are
mutually independent).  For $j\in S_b$:
\begin{equation}\label{eq:aj_decomp}
  \hat{a}_j
  = \frac{1}{n}\X_j^{T}\y
  = \underbrace{\frac{1}{n}\X_j^{T}\X^{(b)}\bal_{S_b}}_{A_j}
    + \underbrace{\frac{1}{n}\X_j^{T}\br^{(b)}}_{B_j}.
\end{equation}
Since $\X_j$ is a column of $\X^{(b)}$, the term $A_j$ depends only
on~$\X^{(b)}$, while $B_j$ depends on $\X^{(b)}$ and on $\br^{(b)}$.

\subsection*{Step 3: Martingale structure}

Consider the filtration
\begin{equation}\label{eq:filtration}
  \mathcal{F}_0=\sigma(\beps,\Z),
  \qquad
  \mathcal{F}_b=\sigma\bigl(\X^{(1)},\dots,\X^{(b)},\beps,\Z\bigr),
  \quad b=1,\dots,B.
\end{equation}
Define the martingale differences
\begin{equation}\label{eq:Db}
  D_b = \E[T_p\mid\mathcal{F}_b]-\E[T_p\mid\mathcal{F}_{b-1}],
  \qquad b=1,\dots,B.
\end{equation}
Then $T_p=\sum_{b=1}^{B}D_b+\E[T_p\mid\mathcal{F}_0]$.  Since
$\E[T_p]=0$ and $\E[T_p\mid\mathcal{F}_B]=T_p$, it follows that
$\E[T_p\mid\mathcal{F}_0]=0$ by iterated expectations (because $\X^{(b)}$ is
independent of $\mathcal{F}_0$ and each $Y_b$ has mean zero) and
$\{M_b=\sum_{b'=1}^{b}D_{b'}\}_{b=0}^{B}$ is a mean-zero martingale
with $M_B=T_p$.

\medskip
\noindent\textbf{Analysis of $D_b$.}\;
The block-diagonal structure ensures that only the terms involving variants
in block~$b$ are affected by $\X^{(b)}$.  Specifically,
\begin{equation}\label{eq:Db_structure}
  D_b
  = \frac{\sqrt{p}}{S_c}\sum_{j\in S_b}c_j
    \Bigl[\hat{a}_j^{2}-\E[\hat{a}_j^{2}\mid\mathcal{F}_{b-1}]\Bigr]
  + \frac{\sqrt{p}}{S_c}\sum_{\substack{j\notin S_b}}c_j
    \Bigl[\E[\hat{a}_j^{2}\mid\mathcal{F}_b]
    -\E[\hat{a}_j^{2}\mid\mathcal{F}_{b-1}]\Bigr].
\end{equation}
For $j\notin S_b$, say $j\in S_{b''}$ with $b''\ne b$, the variable
$\hat{a}_j$ depends on $\X^{(b)}$ only through the ``cross-block'' part of
$\br^{(b'')}$.  However, $\hat{a}_j^{2}$ is a quadratic function
of $\hat{a}_j$, and the dependence on $\X^{(b)}$ enters only through the
linear term $n^{-1}\X_j^{T}\X^{(b)}\bal_{S_b}$.  Because $\X^{(b)}$ is
independent of $\X^{(b'')}$, this dependence contributes at
order~$O(n^{-1})$ per variant, and the total contribution to the second sum
in~\eqref{eq:Db_structure} is of order $O(\sqrt{p}\cdot p\cdot n^{-2}/S_c)
=O(n^{-1})$, which is negligible.

Therefore, to leading order,
\begin{equation}\label{eq:Db_approx}
  D_b \approx Y_b
  = \frac{\sqrt{p}}{S_c}\sum_{j\in S_b}c_j
    \bigl(\hat{a}_j^{2}-\E[\hat{a}_j^{2}]\bigr).
\end{equation}

\subsection*{Step 4: Verification of the Lindeberg condition}

We verify the conditions of the \emph{martingale central limit theorem} (see
Hall \& Heyde, 1980, Theorem~3.2).  We need:
\begin{enumerate}[label=(\alph*)]
\item \emph{Conditional variance convergence:}\;
  $\displaystyle\sum_{b=1}^{B}\E[D_b^{2}\mid\mathcal{F}_{b-1}]
  \xrightarrow{\;P\;}V$.

\item \emph{Lindeberg condition:}\; for every $\delta>0$,
  $\displaystyle\sum_{b=1}^{B}\E[D_b^{2}\,\mathbf{1}(|D_b|>\delta)
  \mid\mathcal{F}_{b-1}]\xrightarrow{\;P\;}0$.
\end{enumerate}

\medskip
\noindent\textbf{Verification of (a).}\;
Using $D_b\approx Y_b$ from~\eqref{eq:Db_approx} and the independence of the
blocks,
\[
  \sum_{b=1}^{B}\E[D_b^{2}\mid\mathcal{F}_{b-1}]
  \approx \sum_{b=1}^{B}\Var(Y_b)
  = \Var(T_p\mid\Z)+\text{cross-block correction}.
\]
The cross-block correction involves $\Cov(Y_b,Y_{b'})$ for $b\ne b'$, which
from Lemma~\ref{lem:cov} is
\[
  \Cov(Y_b,Y_{b'})
  = \frac{p}{S_c^{2}}\sum_{\substack{j\in S_b\\k\in S_{b'}}}c_jc_k
    \,\Cov(\hat{a}_j^{2},\hat{a}_k^{2})
  = \frac{2p}{nS_c^{2}}\sum_{\substack{j\in S_b\\k\in S_{b'}}}c_jc_k
    (a_j^{*}a_k^{*})^{2}+O(n^{-2}).
\]
Together with the within-block terms, the total converges to~$V$ by
Proposition~\ref{prop:variance}.  For the conditional version, the key
observation is that the blocks are of bounded size ($p_b\le p_{\max}$), so
the conditional variance $\E[D_b^{2}\mid\mathcal{F}_{b-1}]$ concentrates
around its unconditional expectation $\E[D_b^{2}]$ at rate
$O(n^{-1/2}p^{-1})$ per block, and the sum of $B=O(p)$ such terms
concentrates at rate $O(n^{-1/2})\to0$.

\medskip
\noindent\textbf{Verification of (b).}\;
Since each block has $p_b\le p_{\max}$ variants and the entries of
$\hat{a}_j$ are sub-Gaussian conditionally on~$\X$, we have the uniform
moment bound
\begin{equation}\label{eq:fourth_moment}
  \E[D_b^{4}\mid\mathcal{F}_{b-1}]
  = O\Bigl(\frac{p^{2}p_b^{2}}{S_c^{4}}\cdot C_b\Bigr)
  = O\bigl(p^{-2}\bigr),
\end{equation}
where $C_b$ is a constant depending on $p_{\max}$, $\bar\lambda$,
$\sigma_\varepsilon^{2}$, and $\max_j|\alpha_j|$ (which is controlled by
Assumption~\ref{ass:effects}).  The bound~\eqref{eq:fourth_moment} follows
because $S_c=\Theta(p^{2})$ in the denominator (recall $S_c=\|\bc\|^{2}$
with $p$ terms).

By the Cauchy--Schwarz and Markov inequalities:
\begin{align}
  \sum_{b=1}^{B}\E[D_b^{2}\,\mathbf{1}(|D_b|>\delta)]
  &\le \sum_{b=1}^{B}\bigl(\E[D_b^{4}]\bigr)^{1/2}
       \bigl(P(|D_b|>\delta)\bigr)^{1/2}
  \notag\\[2pt]
  &\le \sum_{b=1}^{B}O(p^{-1})\cdot\Bigl(\frac{\E[D_b^{4}]}{\delta^{4}}\Bigr)^{1/2}
  \notag\\[2pt]
  &= \sum_{b=1}^{B}O(p^{-1})\cdot O\bigl(p^{-1}\delta^{-2}\bigr)
  = B\cdot O(p^{-2}\delta^{-2})
  = O(p^{-1}\delta^{-2})\to0.
  \label{eq:lindeberg}
\end{align}

\medskip
\noindent\textbf{Conclusion.}\;
Conditions~(a) and~(b) are satisfied.  By the martingale CLT (Hall \& Heyde,
1980, Theorem~3.2),
\[
  T_p = \sum_{b=1}^{B}D_b
  \;\xrightarrow{\;\;d\;\;}\; N(0,V)
  \qquad\text{conditionally on }\Z.
\]

\subsection*{Step 5: Removing the conditioning on \texorpdfstring{$\Z$}{Z}}

We have established that, for almost every realization of $\Z$, the conditional
distribution of $T_p$ given $\Z$ converges weakly to $N(0,V)$.  The
limit~$V$ is the same deterministic constant for all ``typical'' realizations
of~$\Z$ because $c_j\to\ell_j-\bar\ell$ uniformly in~$j$ as
$n_r\to\infty$ (Proposition~\ref{prop:variance}).

Let $\varphi_p(t)=\E[e^{itT_p}]$ denote the characteristic function.  Then
\[
  \varphi_p(t)
  = \E\bigl[\E[e^{itT_p}\mid\Z]\bigr]
  \;\longrightarrow\; e^{-t^{2}V/2}
\]
by the dominated convergence theorem (the integrand is bounded by~$1$, and
$\E[e^{itT_p}\mid\Z]\to e^{-t^{2}V/2}$ in probability).  Hence
$T_p\xrightarrow{d}N(0,V)$ unconditionally.\hfill$\square$

%=============================================================================
\section{Second proof: Chatterjee's second-order Poincar\'e
inequality}\label{sec:chatterjee}
%=============================================================================

We now give a second, independent proof of Theorem~\ref{thm:main} that
provides a quantitative convergence rate.

\subsection{Statement of Chatterjee's inequality}

\begin{theorem}[Chatterjee, 2009, Theorem~2.2]\label{thm:chatterjee}
Let $\bv=(v_1,\dots,v_m)\sim N(\bm{0},\Psi)$ for a positive-definite
$\Psi\in\R^{m\times m}$.  Let $g\in C^{2}(\R^m)$ satisfy
$\E[g(\bv)^{4}]<\infty$, and define
\begin{equation}\label{eq:kappas}
  \kappa_1 = \bigl(\E\bigl[\|\nabla g(\bv)\|^{4}\bigr]\bigr)^{1/4},
  \qquad
  \kappa_2 = \bigl(\E\bigl[\|\nabla^{2}g(\bv)\|_{\op}^{4}\bigr]\bigr)^{1/4},
  \qquad
  \psi^{2} = \Var\{g(\bv)\}.
\end{equation}
If $w\sim N\bigl(\E[g(\bv)],\,\psi^{2}\bigr)$, then
\begin{equation}\label{eq:chatterjee_bound}
  \dTV\bigl\{g(\bv),\,w\bigr\}
  \le 2\sqrt{5}\;\frac{\|\Psi\|_{\op}^{3/2}\,\kappa_1\,\kappa_2}{\psi^{2}}.
\end{equation}
\end{theorem}

\subsection{Setup}\label{ssec:setup_chatterjee}

Condition on $\Z$ and set
\begin{equation}\label{eq:vdef}
  \bv = \bigl(\mathrm{vec}(\X^{T}),\,\beps\bigr)\in\R^{m},
  \qquad m=(p+1)n.
\end{equation}
Then $\bv\sim N(\bm{0},\Psi)$ with
$\Psi=I_n\otimes\diag(\bSig,\sigma_\varepsilon^{2})$, so that
\begin{equation}\label{eq:Psi_op}
  \|\Psi\|_{\op}
  = \max\bigl(\|\bSig\|_{\op},\,\sigma_\varepsilon^{2}\bigr)
  \le \bar\lambda\vee\sigma_\varepsilon^{2}
  =: \Lambda.
\end{equation}
Define $g:\R^{m}\to\R$ by
\begin{equation}\label{eq:gdef}
  g(\bv)
  = g(\X,\beps)
  = \frac{1}{n^{2}S_c}\,\y^{T}\X D\X^{T}\y
  = \frac{1}{n^{2}S_c}\,(\X\bal+\beps)^{T}\X D\X^{T}(\X\bal+\beps).
\end{equation}
This is a polynomial of degree~4 in the entries of~$\bv$, hence
$g\in C^{\infty}(\R^m)$ and all moments are finite.

\subsection{Gradient bound ($\kappa_1$)}\label{ssec:kappa1}

We compute the gradient of~$g$ with respect to the entries of~$\bv$.  Write
$\X=(x_{ij})_{\substack{1\le i\le n\\1\le j\le p}}$ and
$\beps=(\varepsilon_1,\dots,\varepsilon_n)^{T}$, and recall $\y=\X\bal+\beps$.

\medskip
\noindent\textbf{Partial derivatives with respect to $x_{ij}$.}\;
Let $E_{ij}\in\R^{n\times p}$ be the matrix with $(E_{ij})_{i'j'}=\delta_{ii'}\delta_{jj'}$.
Since
\begin{equation}\label{eq:partials_x1}
  \frac{\partial}{\partial x_{ij}}\y^{T}\X D\X^{T}\y
  = 2\bal^{T}E_{ij}^{T}D\X^{T}\y\cdot y_i
    + 2y_i\,(\X D)_{ij}\,\bx_i^{T}\bal
    + 2y_i\be_j^{T}D\X^{T}\y\cdot\alpha_j
    \quad\text{(schematically)},
\end{equation}
more precisely, using $\frac{\partial}{\partial x_{ij}}\|\X^{T}\y\|_{D}^{2}
= 2y_i(\X D\X^{T}\y)_i$ (differentiation of the quadratic form
$\y^{T}\X D\X^{T}\y$ with respect to the matrix entries), and the fact that
$\y$ depends on $x_{ij}$ through $\X\bal$, the full derivative is
\begin{equation}\label{eq:dg_dxij}
  \frac{\partial g}{\partial x_{ij}}
  = \frac{2}{n^{2}S_c}
    \Bigl[\alpha_j\,\be_i^{T}\X D\X^{T}\y
    + y_i\,\be_j^{T}D\X^{T}\y
    + y_i\,\be_j^{T}D\X^{T}\X\bal\cdot\delta_{jj'}\text{ terms}\Bigr].
\end{equation}
Following the analysis of equations~(S5)--(S8) in the Supplement of
Dicker~(2014)---who performs an identical gradient calculation for the
function $n^{-2}\|\X^{T}\y\|^{2}$ and obtains
$\kappa_1=O(n^{-1/2})$ times polynomial factors in the spectral norm
$\lambda_1=\|n^{-1}\X^{T}\X\|_{\op}$ and the signal/noise parameters---we
replace Dicker's identity matrix with $D/S_c$ and obtain
\begin{align}
  \|\nabla_{\X}g\|^{2}
  &\le \frac{C}{n^{4}S_c^{2}}
    \Bigl[\|D\|_{\op}^{2}\bigl(\|\y\|^{2}\|\X\|_{\op}^{2}
    + \tau^{2}\|\X\|_{\op}^{4}\bigr)
    \|\y\|^{2}\Bigr],
  \label{eq:gradX_bound}
\end{align}
for an absolute constant~$C>0$.

\medskip
\noindent\textbf{Partial derivatives with respect to $\varepsilon_k$.}\;
Since $\partial\y/\partial\varepsilon_k=\be_k$ (the $k$-th standard basis
vector in~$\R^n$), we have
\begin{equation}\label{eq:dg_deps}
  \frac{\partial g}{\partial\varepsilon_k}
  = \frac{2}{n^{2}S_c}\,\be_k^{T}\X D\X^{T}\y,
\end{equation}
whence
\begin{equation}\label{eq:gradeps_bound}
  \sum_{k=1}^{n}\Bigl(\frac{\partial g}{\partial\varepsilon_k}\Bigr)^{\!2}
  = \frac{4}{n^{4}S_c^{2}}\,\|\X D\X^{T}\y\|^{2}
  \le \frac{4\|D\|_{\op}^{2}}{n^{4}S_c^{2}}
    \,\|\X\|_{\op}^{4}\,\|\y\|^{2}.
\end{equation}

\medskip
\noindent\textbf{Combined gradient bound.}\;
Combining~\eqref{eq:gradX_bound} and~\eqref{eq:gradeps_bound}, and using
$\|\y\|^{2}=O_P(n\sigma_y^{2})$,
$\|\X\|_{\op}^{2}=O_P(n\bar\lambda)$, $\|D\|_{\op}=\max_j|c_j|=O(1)$,
and $S_c=\Theta(p)$, we obtain
\begin{equation}\label{eq:grad_norm_bound}
  \|\nabla g(\bv)\|^{2}
  = O_P\Bigl(\frac{1}{n\,p^{2}}\Bigr).
\end{equation}
Therefore, taking fourth moments:
\begin{equation}\label{eq:kappa1_bound}
  \boxed{\;\kappa_1
  = \bigl(\E[\|\nabla g(\bv)\|^{4}]\bigr)^{1/4}
  = O\bigl(n^{-1/2}p^{-1}\bigr).\;}
\end{equation}
The existence and finiteness of $\E[\|\nabla g(\bv)\|^{4}]$ is guaranteed by
Assumption~\ref{ass:spectral} and the bounded block-size condition
(Assumption~\ref{ass:block}), which imply that $\lambda_1$ has bounded moments
of all orders.  Indeed, $\|\X\|_{\op}^{2}=\max_b\|\X^{(b)}\|_{\op}^{2}$ and
each $\|\X^{(b)}\|_{\op}^{2}$ is the largest eigenvalue of a
Wishart$(n,\bSig_b)$ matrix of fixed dimension $p_b\le p_{\max}$; by the
Bai--Yin law and exponential tail bounds for the largest eigenvalue
(see, e.g., Davidson \& Szarek, 2001), all polynomial moments of
$\|\X\|_{\op}^{2}$ are uniformly bounded.

\subsection{Hessian bound ($\kappa_2$)}\label{ssec:kappa2}

Following the analysis of equations~(S9)--(S12) of Dicker~(2014), we bound
the operator norm of the Hessian $\nabla^{2}g(\bv)$.  Let
$\tilde U=(\tilde u\;\;U)\in\R^{n\times(p+1)}$ be a matrix with
$\|\tilde U\|_F=1$, where $\tilde u\in\R^{n}$ and $U\in\R^{n\times p}$.
Define the directional differential operator
$D_{\tilde U}=\sum_{i,j}U_{ij}\frac{\partial}{\partial x_{ij}}
+\sum_k\tilde u_k\frac{\partial}{\partial\varepsilon_k}$.
Then
\begin{equation}\label{eq:Hessian_op}
  \|\nabla^{2}g(\bv)\|_{\op}
  = \sup_{\|\tilde U\|_F=1}|D_{\tilde U}^{2}\,g(\bv)|.
\end{equation}

The function $g$ is a degree-4 polynomial in the entries of~$\bv$, so
$D_{\tilde U}^{2}g$ is a degree-2 polynomial (a quadratic form in $\bv$
plus a constant).  The same computation as in Dicker~(2014), equations
(S9)--(S11), adapted to our weight matrix $D/S_c$, gives
\begin{align}
  |D_{\tilde U}^{2}g(\bv)|
  &\le \frac{C'}{n^{2}S_c}
    \Bigl[(\lambda_1+1)(\tau+\sigma_\varepsilon)^{2}
    + \frac{\normsq{\beps}}{n}\Bigr]
  \label{eq:DU2g_bound}
\end{align}
for a constant $C'$ depending on $\|D\|_{\op}$ (which is $O(1)$).  Since this
bound is uniform over $\tilde U$ with $\|\tilde U\|_F=1$:
\begin{equation}\label{eq:Hessian_norm_bound}
  \|\nabla^{2}g(\bv)\|_{\op}
  = O_P\Bigl(\frac{1}{n\,p}\Bigr).
\end{equation}
Taking fourth moments (using the same moment bounds as for~$\kappa_1$):
\begin{equation}\label{eq:kappa2_bound}
  \boxed{\;\kappa_2
  = \bigl(\E[\|\nabla^{2}g(\bv)\|_{\op}^{4}]\bigr)^{1/4}
  = O\bigl(n^{-1}p^{-1}\bigr).\;}
\end{equation}

\subsection{Variance lower bound}\label{ssec:psi}

From Proposition~\ref{prop:variance},
\begin{equation}\label{eq:psi2_bound}
  \psi^{2}
  = \Var(g(\bv))
  = \Var(\hat\sigma_\alpha^{2}\mid\Z)
  = \frac{V}{p}+o\Bigl(\frac{1}{p}\Bigr)
  = \Theta(p^{-1}).
\end{equation}

\subsection{Applying Theorem~\ref{thm:chatterjee}}

Substituting~\eqref{eq:Psi_op},~\eqref{eq:kappa1_bound},~\eqref{eq:kappa2_bound},
and~\eqref{eq:psi2_bound} into~\eqref{eq:chatterjee_bound}:
\begin{align}
  \dTV\bigl(\hat\sigma_\alpha^{2},\,w\bigr)
  &\le 2\sqrt{5}\;\frac{\Lambda^{3/2}\,\kappa_1\,\kappa_2}{\psi^{2}}
  \notag\\[4pt]
  &= O(1)\cdot\frac{(n^{-1/2}p^{-1})(n^{-1}p^{-1})}{p^{-1}}
  \notag\\[4pt]
  &= O\bigl(n^{-3/2}p^{-1}\bigr)
  \;\longrightarrow\; 0,
  \label{eq:dTV_rate}
\end{align}
where $w\sim N(\sigma_{\alpha,0}^{2},\psi^{2})$.

Since $\dTV\to0$ implies convergence in distribution, and since
$\psi^{2}=V/p+o(1/p)$, we conclude that
\begin{equation}\label{eq:standardised}
  \frac{\hat\sigma_\alpha^{2}-\sigma_{\alpha,0}^{2}}
       {\sqrt{\Var(\hat\sigma_\alpha^{2}\mid\Z)}}
  \;\xrightarrow{\;\;d\;\;}\; N(0,1).
\end{equation}

\subsection{Removing the conditioning on \texorpdfstring{$\Z$}{Z}}

The same argument as in Step~5 of Section~\ref{sec:mCLT} applies.  The
limit variance~$V$ is a deterministic constant that does not depend on the
realization of~$\Z$, and the characteristic function argument shows that
\[
  \sqrt{p}\,\bigl(\hat\sigma_\alpha^{2}-\sigma_{\alpha,0}^{2}\bigr)
  \;\xrightarrow{\;\;d\;\;}\; N(0,V)
\]
unconditionally.  Moreover, the bound~\eqref{eq:dTV_rate} gives the
\emph{quantitative rate}
\begin{equation}\label{eq:rate}
  \dTV\bigl(\sqrt{p}(\hat\sigma_\alpha^{2}-\sigma_{\alpha,0}^{2}),\;
  N(0,V)\bigr)
  = O\bigl(n^{-3/2}p^{-1}\bigr)
  = O\bigl(n^{-5/2}\bigr)
\end{equation}
when $p=\Theta(n)$.  This completes the second proof of
Theorem~\ref{thm:main}.\hfill$\square$

%=============================================================================
\section{Explicit form of the asymptotic variance}\label{sec:Vexplicit}
%=============================================================================

Under the block-diagonal structure~\eqref{eq:blockSig}, $\Sigma_{jk}=0$
whenever $j$ and $k$ belong to different blocks.  Hence the double sum
in~\eqref{eq:Vdef} splits into within-block and cross-block contributions:
\begin{align}
  V &= \lim_{p\to\infty}\;\frac{2\gamma}{V_\ell^{2}}\cdot
    \frac{1}{p^{2}}\Biggl[
    \sum_{b=1}^{B}\sum_{\substack{j\in S_b\\k\in S_b}}
      (\ell_j-\bar\ell)(\ell_k-\bar\ell)
      \Bigl(\frac{\Sigma_{jk}\sigma_y^{2}}{n}+a_j^{*}a_k^{*}\Bigr)^{\!2}
    \notag\\[6pt]
    &\qquad\qquad\qquad
    + \sum_{\substack{b,b'=1\\b\ne b'}}^{B}
      \sum_{\substack{j\in S_b\\k\in S_{b'}}}
        (\ell_j-\bar\ell)(\ell_k-\bar\ell)\,(a_j^{*})^{2}(a_k^{*})^{2}
    \Biggr],
  \label{eq:Vexplicit}
\end{align}
where $\gamma=\lim p/n$.

\medskip\noindent\textbf{Within-block terms.}\;
For $j,k\in S_b$, the covariance $\Sigma_{jk}$ is potentially non-zero.  The
leading-order contribution from block~$b$ involves the fourth moments of the
Wishart matrix $W_b=(\X^{(b)})^{T}\X^{(b)}\sim\text{Wishart}_{p_b}(n,\bSig_b)$,
which are computed exactly by Proposition~S1 of Dicker~(2014).  Since
$p_b\le p_{\max}$, these are $O(1)$ quantities.

\medskip\noindent\textbf{Cross-block terms.}\;
For $b\ne b'$, $\Sigma_{jk}=0$ and only the product $(a_j^{*})^{2}(a_k^{*})^{2}$
survives.  This factorises as
\[
  \sum_{\substack{b\ne b'}}\sum_{\substack{j\in S_b\\k\in S_{b'}}}
  (\ell_j-\bar\ell)(\ell_k-\bar\ell)(a_j^{*})^{2}(a_k^{*})^{2}
  = \Bigl(\sum_{j=1}^{p}(\ell_j-\bar\ell)(a_j^{*})^{2}\Bigr)^{\!2}
    - \sum_{b=1}^{B}\Bigl(\sum_{j\in S_b}(\ell_j-\bar\ell)(a_j^{*})^{2}\Bigr)^{\!2},
\]
which simplifies the computation considerably.

\begin{remark}[Identifiability]\label{rem:identity}
When $\bSig=I_p$, we have $\ell_j=1$ for all~$j$ and $V_\ell=0$, so
Assumption~\ref{ass:ldvar} fails.  This reflects the well-known fact that LDSC
requires heterogeneity in LD scores across variants for identifiability.  The
block-diagonal structure with non-trivial within-block LD (i.e., $\bSig_b$
having non-degenerate eigenvalues) ensures that the $\ell_j$ are
heterogeneous, which is the empirically relevant regime.
\end{remark}

\begin{remark}[Comparison with random-effects models]\label{rem:random}
Under a standard random-effects model where
$\alpha_j\overset{\mathrm{iid}}{\sim}(0,\sigma_\alpha^{2}/p)$, the
pseudo-true value satisfies $\sigma_{\alpha,0}^{2}=\sigma_\alpha^{2}+O(n^{-1})$,
so the LDSC estimator is consistent for the variance component.  Under our
fixed-effects framework, $\sigma_{\alpha,0}^{2}$ depends on the specific
alignment between the effect vector~$\bal$ and the LD structure, and its
interpretation must be qualified accordingly.
\end{remark}

%=============================================================================
\section*{References}
%=============================================================================

\begin{description}[leftmargin=0pt,itemsep=6pt]
\item[] Bai, Z.\ \& Yin, Y.\ Q.\ (1993). Limit of the smallest eigenvalue of
  a large dimensional sample covariance matrix.
  \textit{Ann.\ Probab.}\ \textbf{21}, 1275--1294.

\item[] Bulik-Sullivan, B.\ K.\ et~al.\ (2015). LD score regression
  distinguishes confounding from polygenicity in genome-wide association
  studies. \textit{Nature Genetics}\ \textbf{47}, 291--295.

\item[] Chatterjee, S.\ (2009). Fluctuations of eigenvalues and second order
  Poincar\'e inequalities. \textit{Probab.\ Theory Relat.\ Fields}\
  \textbf{143}, 1--40.

\item[] Davidson, K.\ R.\ \& Szarek, S.\ J.\ (2001). Local operator theory,
  random matrices and Banach spaces. In \textit{Handbook of the Geometry of
  Banach Spaces}, Vol.~1, pp.~317--366. Elsevier.

\item[] Dicker, L.\ H.\ (2014). Variance estimation in high-dimensional
  linear models. \textit{Biometrika}\ \textbf{101}, 269--284.

\item[] Hall, P.\ \& Heyde, C.\ C.\ (1980). \textit{Martingale Limit Theory
  and Its Application}. Academic Press, New York.

\item[] Letac, G.\ \& Massam, H.\ (2004). All invariant moments of the
  Wishart distribution. \textit{Scand.\ J.\ Statist.}\ \textbf{31}, 295--318.
\end{description}

\end{document}
